\section{The isotropic random-vector theorem}
\label{sec:mss}

This appendix proves the precise form of the Marcus--Spielman--Srivastava
theorem used in Section~\ref{sec:sparsification}.  The proof proceeds
through expected characteristic polynomials, an elementary interlacing
argument, and a barrier estimate.  In particular, the barrier estimate
below is proved directly from partial fractions; no representation
theorem for bivariate stable polynomials is required.

\begin{theorem}[Marcus--Spielman--Srivastava, {\cite[Theorem~1.4]{MSS}}]
\label{thm:mss}
Let $\epsilon>0$, and let $v_1,\ldots,v_m$ be independent random vectors
with finite support in $\mathbb C^d$.  If
\begin{equation}
 \sum_{i=1}^m\mathbb E[v_iv_i^*]=I_d,
 \qquad \mathbb E\|v_i\|^2\leq\epsilon\quad(1\leq i\leq m),
 \label{eq:mss-hypotheses}
\end{equation}
then
\begin{equation}
 \mathbb P\left\{
 \left\|\sum_{i=1}^m v_iv_i^*\right\|
       \leq(1+\sqrt\epsilon)^2\right\}>0.
 \label{eq:mss-conclusion}
\end{equation}
Here $*$ denotes conjugate transpose and the norm is the operator norm.
In particular, some outcome satisfies the corresponding
positive-semidefinite upper bound.  Real vectors are a special case.
\end{theorem}

The case $d=0$ is immediate, so the proof below assumes $d\geq1$.
We will use polynomial factorization, elementary compactness, and
standard finite-dimensional linear algebra.  The polynomial and
analytic lemmas needed beyond those facts are proved below.

\subsection{Polynomial limits and stability}

Write $\mathbb H=\{z\in\mathbb C:\operatorname{Im}z>0\}$.
A nonzero polynomial $p$ in several variables is \emph{stable} if it
does not vanish when all its arguments lie in $\mathbb H$.  It is
\emph{real stable} if, in addition, its coefficients are real.
A nonzero real univariate polynomial is stable exactly when all its
roots are real: nonreal roots of a real polynomial occur in conjugate
pairs.

\begin{lemma}[Limits of univariate polynomials]
\label{lem:mss-polynomial-limits}
Suppose nonzero univariate polynomials $f_\nu$ have uniformly bounded
degrees and converge coefficientwise to a nonzero polynomial $f$.
If all $f_\nu$ have no zeros in $\mathbb H$, neither does $f$.
If all $f_\nu$ are real-rooted, so is $f$.
For monic real-rooted polynomials of a fixed degree, their ordered
roots depend continuously on their coefficients.
\end{lemma}

\begin{proof}
For the first assertion choose $a\in\mathbb H$ with $f(a)\ne0$.
Pass to a subsequence on which the degree is constant, and enumerate
its roots $\lambda_{\nu,1},\ldots,\lambda_{\nu,e}$ with multiplicity.
By compactness of the complex plane together with the point at
infinity, a further subsequence makes every root converge, allowing
infinite limits.  Every finite limit lies outside $\mathbb H$.
In the factorization
\[
 \frac{f_\nu(z)}{f_\nu(a)}
   =\prod_{k=1}^e\frac{z-\lambda_{\nu,k}}{a-\lambda_{\nu,k}},
\]
a root tending to infinity contributes a factor tending to $1$, and
a root tending to a finite $\lambda_k$ contributes
$(z-\lambda_k)/(a-\lambda_k)$.  The limiting identity factors
$f(z)/f(a)$ with all its roots outside $\mathbb H$.
For real-rooted $f_\nu$, the same proof uses the extended real line;
all finite limiting roots are then real.

For the last assertion, a bounded set of monic coefficient vectors
has bounded roots.  Indeed a root of
$z^d+\sum_{k<d}a_kz^k$ has absolute value at most
$1+\max_k|a_k|$, as is seen by summing a geometric series when the
opposite strict inequality holds.  Any convergent subsequence of
ordered root tuples factors the limiting polynomial.  Its limiting
tuple must therefore be the unique ordered root tuple of that
polynomial.  This proves continuity.
\end{proof}

\begin{lemma}[Elementary stability operations]
\label{lem:mss-stability}
The following assertions hold.
\begin{enumerate}
\item If $A_1,\ldots,A_m$ are Hermitian positive-semidefinite matrices
and $C$ is Hermitian, then
\[
 \det\left(xI+C+\sum_{i=1}^m z_iA_i\right)
\]
is real stable in $x,z_1,\ldots,z_m$.
If $\sum_iA_i$ is positive definite, then
$\det(\sum_i z_iA_i)$ is also real stable.
\item If $p$ is real stable, then $(1-\partial_i)p$ is real stable.
\item A nonzero univariate polynomial obtained by setting all but one
variable of a real stable polynomial to real constants is real-rooted.
\end{enumerate}
\end{lemma}

\begin{proof}
For the first assertion, when all variables are in $\mathbb H$,
the imaginary part of the matrix is
\[
 (\operatorname{Im}x)I+\sum_i(\operatorname{Im}z_i)A_i\succ0.
\]
Such a matrix is nonsingular: a nonzero kernel vector would have a
strictly positive imaginary part in its quadratic form, a
contradiction.  The same argument applies to the second determinant
because $\sum_iA_i\succ0$.  For real arguments the matrices are
Hermitian and their determinants are real; consequently the
polynomial coefficients are real.

For the second assertion, fix all arguments of $p$ in $\mathbb H$,
and regard $p$ as a polynomial in argument $i$ alone.  This slice is
nonzero, and its roots $\alpha_k$ have nonpositive imaginary parts.
If it is nonconstant, factorization gives
\begin{equation}
 \frac{\partial_i p}{p}
   =\sum_k\frac1{z_i-\alpha_k},
 \qquad
 \operatorname{Im}\frac{\partial_i p}{p}<0.
 \label{eq:mss-log-derivative-upper}
\end{equation}
If the slice is constant, the quotient is zero.  In either case it
cannot equal $1$, so $(1-\partial_i)p$ is nonzero at this point.
The operation preserves real coefficients.

For the last assertion, replace each fixed real constant $a$ by
$a+\mathrm i\eta$, $\eta>0$.  The resulting univariate polynomials
have no zeros in $\mathbb H$ and converge coefficientwise to the
stated nonzero real polynomial.  Apply
Lemma~\ref{lem:mss-polynomial-limits}.
\end{proof}

\subsection{Expected characteristic polynomials}

For positive-semidefinite Hermitian $d$-by-$d$ matrices
$A_1,\ldots,A_m$, define
\begin{equation}
 \mu[A_1,\ldots,A_m](x)
 =\left.
   \left(\prod_{i=1}^m(1-\partial_{z_i})\right)
   \det\left(xI_d+\sum_{i=1}^m z_iA_i\right)
   \right|_{z_1=\cdots=z_m=0}.
 \label{eq:mss-mixed-polynomial}
\end{equation}

\begin{lemma}[The expectation identity and real roots]
\label{lem:mss-expectation}
The polynomial in \eqref{eq:mss-mixed-polynomial} is monic of degree
$d$ and real-rooted.  If $v_1,\ldots,v_m$ are independent finitely
supported random vectors and $A_i=\mathbb E[v_iv_i^*]$, then
\begin{equation}
 \mathbb E\det\left(xI_d-\sum_{i=1}^m v_iv_i^*\right)
 =\mu[A_1,\ldots,A_m](x).
 \label{eq:mss-expectation-identity}
\end{equation}
\end{lemma}

\begin{proof}
The determinant in \eqref{eq:mss-mixed-polynomial} is stable by
Lemma~\ref{lem:mss-stability}; applying all the operators preserves
stability.  Specializing all $z_i$ to zero leaves a nonzero real-rooted
polynomial.  To see directly both that it is nonzero and that it is
monic of degree $d$, the coefficient of $x^d$ in the determinant is
$1$, independent of the $z_i$, and every term involving at least one
$z_i$ has smaller $x$-degree.  The operations leave that leading
coefficient unchanged.  The specialization can be performed all at
once by substituting $z_i=\mathrm i\eta$ and using
Lemma~\ref{lem:mss-polynomial-limits}.

For the identity, let $M$ be any square matrix, let $v$ be a random
vector, and put $A=\mathbb E[vv^*]$.  The rank-one determinant identity
and the derivative formula for the determinant give
\begin{align*}
 \mathbb E\det(M-vv^*)
 &=\det M-\operatorname{tr}(\operatorname{adj}(M)A)\\
 &=\left.(1-\partial_z)\det(M+zA)\right|_{z=0}.
\end{align*}
Here $\operatorname{adj}(M)$ is the adjugate.  These identities hold
also for singular $M$: they are polynomial identities in its entries,
or follow directly by multilinearity of the determinant in its
columns.  Apply this formula conditionally to $v_m$, then $v_{m-1}$,
and so forth.  Independence keeps each covariance $A_i$ unchanged
under conditioning.  Finite sums, differentiation, and the
specializations commute, giving
\eqref{eq:mss-expectation-identity}.
\end{proof}

\subsection{Interlacing and the choice of an outcome}

For a monic real-rooted degree-$d$ polynomial $f$, write its roots as
$\lambda_1(f)\leq\cdots\leq\lambda_d(f)$, counting multiplicity.
A finite collection of such polynomials has a \emph{common
interlacing} if there are real numbers
$c_1\leq\cdots\leq c_{d-1}$ such that
\[
 \lambda_k(f)\leq c_k\leq\lambda_{k+1}(f)
 \quad(1\leq k<d)
\]
for every polynomial in the collection.

\begin{lemma}[An elementary interlacing criterion]
\label{lem:mss-common-interlacing}
If every convex combination of finitely many monic degree-$d$
polynomials is real-rooted, they have a common interlacing.
Consequently, for any probability weights on this collection, some
polynomial of positive weight has largest root at most the largest
root of the weighted average.
\end{lemma}

\begin{proof}
First consider two members $p,q$ and
$r_t=(1-t)p+tq$, $0\leq t\leq1$.
Fix a real $x$ that is not a root of either $p$ or $q$, and let $N(t)$
count the roots of $r_t$ strictly below $x$, with multiplicity.
By ordered-root continuity, $N$ is locally constant wherever
$r_t(x)\ne0$.  The affine function $r_t(x)$ vanishes at most once.
If it vanishes at $t_0\in(0,1)$, then
$b=q(x)-p(x)\ne0$.

We claim that $x$ is a simple root of $r_{t_0}$.  Otherwise let its
multiplicity be $\ell\geq2$ and write
\[
 r_{t_0}(x+w)=aw^\ell+O(w^{\ell+1}),\qquad a\ne0.
\]
For $h$ tending to zero with fixed sign $\sigma\in\{-1,1\}$,
\[
 |h|^{-1}r_{t_0+h}\bigl(x+|h|^{1/\ell}z\bigr)
 \longrightarrow az^\ell+\sigma b
\]
coefficientwise.  Every polynomial on the left is real-rooted; hence
so is its nonzero limit by Lemma~\ref{lem:mss-polynomial-limits}.
For $\ell=2$, choose the sign so that $\sigma b/a>0$, obtaining
two nonreal roots.  For $\ell\geq3$, the equation
$z^\ell=-\sigma b/a\ne0$ has nonreal roots for either sign.
This contradiction proves the claim.

Thus $N$ can change by at most one, at that single simple crossing.
It follows that $|N(0)-N(1)|\leq1$.
If $\lambda_k(p)>\lambda_{k+1}(q)$ for some $k<d$, one can choose
$x$ strictly between these roots and avoiding all endpoint roots.
Then $N(0)\leq k-1$ and $N(1)\geq k+1$, a contradiction.
Applying this reasoning to every pair gives
\[
 \lambda_k(p_i)\leq\lambda_{k+1}(p_j)
 \quad\hbox{for all }i,j\hbox{ and }k<d.
\]
Therefore $c_k=\max_i\lambda_k(p_i)$ is a common interlacing.
This proof includes repeated roots because all interlacing
inequalities are weak.

For the last assertion discard zero-weight polynomials and denote
the weighted average by $P$, with largest root $\rho$.
If $d=1$, $\rho$ is the weighted average of the individual roots,
which proves the assertion.  Otherwise put $c=c_{d-1}$.
Every $p_i(c)\leq0$, since its first $d-1$ roots are at most $c$
and its last root is at least $c$.  Thus $P(c)\leq0$, and, since
$P$ is monic, $\rho\geq c$.
For each $x>\rho$, $P(x)>0$, so at least one $p_i(x)>0$.
Its first $d-1$ roots are below $x$; positivity therefore implies
$\lambda_d(p_i)<x$.  Taking the minimum over the finite collection
and letting $x\downarrow\rho$ proves
$\min_i\lambda_d(p_i)\leq\rho$.
\end{proof}

\begin{lemma}[Selection from independent supports]
\label{lem:mss-selection}
For independent finitely supported vectors $v_1,\ldots,v_m$, some
outcome of positive probability has largest eigenvalue of
$\sum_i v_iv_i^*$ at most the largest root of
$\mu[\mathbb E v_1v_1^*,\ldots,\mathbb E v_mv_m^*]$.
\end{lemma}

\begin{proof}
Delete zero-probability support points.  For a fixed choice of the
first $k$ vectors, consider the expected characteristic polynomial
over the remaining vectors.  It is monic of degree $d$ and
real-rooted by Lemma~\ref{lem:mss-expectation}, treating the first
$k$ vectors as deterministic.  Its children are the conditional
polynomials formed by fixing one further vector to each of its
possible values.

Every convex combination of these children is again real-rooted:
it is the expected characteristic polynomial for the same family,
with the next vector assigned the probabilities of that convex
combination.  Lemma~\ref{lem:mss-common-interlacing} therefore lets
us choose a positive-probability child whose largest root does not
increase.  Repeating for $k=0,\ldots,m-1$ yields a complete outcome.
The probability of that outcome is a product of finitely many
positive support probabilities, so is positive.  At the leaf, the
polynomial is the characteristic polynomial of
$\sum_i v_iv_i^*$ and its largest root is the largest eigenvalue.
\end{proof}

\subsection{Barrier functions: a direct proof}

For a real stable polynomial $p$, let
\[
 \operatorname{Ab}_p
 =\{z\in\mathbb R^m:p(z+s)>0
                    \text{ for every }s\in[0,\infty)^m\}.
\]
Such a point is said to be \emph{above the roots}.  At a point of
this set define
\[
 \Phi_p^i(z)=\frac{\partial_i p(z)}{p(z)}.
\]
The set $\operatorname{Ab}_p$ is closed under adding a vector with
nonnegative coordinates.  The derivatives below are well defined
locally because $p$ is positive at all the points under discussion.

\begin{lemma}[Positive partial fractions]
\label{lem:mss-partial-fractions}
Let $p$ be real stable, $z\in\operatorname{Ab}_p$, and let $i,j$ be
coordinate indices.  As a rational function of $t$,
\begin{equation}
 \Phi_p^i(z+te_j)
 =c+\sum_{\ell=1}^s\frac{c_\ell}{t-\lambda_\ell},
 \qquad c\geq0,\quad c_\ell>0,\quad\lambda_\ell<0,
 \label{eq:mss-positive-partial-fractions}
\end{equation}
where the sum may be empty.  In particular, for $t\geq0$ this
function is nonnegative, nonincreasing, and convex.
\end{lemma}

\begin{proof}
Set $h(t)=p(z+te_j)$ and $a(t)=\partial_i p(z+te_j)$.
The polynomial $h$ is real-rooted by
Lemma~\ref{lem:mss-stability}, and is nonzero since $h(0)>0$.
It is positive on $[0,\infty)$, so every root is strictly negative.
Thus the real rational function $R=a/h$ has only negative real
poles, after cancelling common factors.

We next show that $\operatorname{Im}R(t)\leq0$ for $t\in\mathbb H$.
In \eqref{eq:mss-log-derivative-upper}, use the point with argument
$z_j+t$ in coordinate $j$ and argument $z_k+\mathrm i\eta$ in every
other coordinate.  All coordinates have positive imaginary parts,
and the logarithmic derivative in direction $i$ has nonpositive
imaginary part.  Let $\eta\downarrow0$.
Since $h(t)\ne0$ for $t\in\mathbb H$, the quotient converges to
$R(t)$ and proves the desired inequality.

Also $R(t)\geq0$ for real $t\geq0$.  Indeed
$g(s)=p(z+te_j+se_i)$ is a nonzero real-rooted polynomial, positive
for $s\geq0$.  Factoring it with its negative real roots $\beta_k$
gives
\[
 R(t)=\frac{g'(0)}{g(0)}=\sum_k\frac1{-\beta_k}\geq0;
\]
the empty sum covers a constant $g$.

It remains to classify this rational function, which can be done
directly.  At a pole $\lambda$ of order $b\geq1$, let its highest
principal part be $a_0(t-\lambda)^{-b}$, with $a_0\ne0$ real.
At $t=\lambda+r\exp(\mathrm i\theta)$, as $r\downarrow0$, the
leading imaginary part is
$-a_0r^{-b}\sin(b\theta)$.
For $b\geq2$ this has both signs for $0<\theta<\pi$, contradicting
$\operatorname{Im}R\leq0$.  Hence every pole is simple, and the same
inequality for $b=1$ forces its residue to be positive.

Subtract these simple principal parts.  The remainder is a real
polynomial.  If its degree were at least two, its leading term along
large rays $t=r\exp(\mathrm i\theta)$ would again have imaginary
parts of both signs; the subtracted fractions are $O(1/r)$ and
cannot affect this conclusion.  Thus the remainder has degree at
most one, and a linear remainder must have nonpositive slope.
The inequality $R(t)\geq0$ for all real $t\geq0$ rules out a
negative slope and forces its constant term to be nonnegative.
This proves \eqref{eq:mss-positive-partial-fractions}.

Finally, direct differentiation on $t\geq0$ gives
\[
 R'(t)=-\sum_\ell\frac{c_\ell}{(t-\lambda_\ell)^2}\leq0,
 \qquad
 R''(t)=2\sum_\ell\frac{c_\ell}{(t-\lambda_\ell)^3}\geq0.
\]
Together with $R(t)\geq0$, these are the asserted properties.
\end{proof}

\begin{lemma}[One barrier shift]
\label{lem:mss-barrier-shift}
Let $p$ be real stable, let $z\in\operatorname{Ab}_p$, and suppose
$\delta>0$ satisfies
\begin{equation}
 \Phi_p^j(z)+\frac1\delta\leq1.
 \label{eq:mss-shift-hypothesis}
\end{equation}
For $q=(1-\partial_j)p$ and $w=z+\delta e_j$, one has
$w\in\operatorname{Ab}_q$ and
\begin{equation}
 \Phi_q^i(w)\leq\Phi_p^i(z)\qquad\text{for every }i.
 \label{eq:mss-shift-conclusion}
\end{equation}
\end{lemma}

\begin{proof}
The polynomial $q$ is real stable by
Lemma~\ref{lem:mss-stability}.  Applying the monotonicity in
Lemma~\ref{lem:mss-partial-fractions} successively in each coordinate,
for every vector $s\geq0$ we have
\[
 \Phi_p^j(z+s)\leq\Phi_p^j(z)\leq1-1/\delta<1.
\]
It follows that
$q(z+s)=p(z+s)(1-\Phi_p^j(z+s))>0$.
Thus $z$, and therefore also $w$, belongs to $\operatorname{Ab}_q$.

Fix $i$ and put $f(t)=\Phi_p^i(z+te_j)$.
Commuting the mixed derivatives of $\log p$ in a neighborhood of
$w$ gives
\[
 \Phi_q^i(w)
 =f(\delta)-\frac{f'(\delta)}{1-\Phi_p^j(w)}.
\]
By monotonicity, the denominator is at least $1/\delta$, and
$f'(\delta)\leq0$.  By convexity,
$f(0)\geq f(\delta)-\delta f'(\delta)$; this also follows by
integrating the nondecreasing derivative on $[0,\delta]$.
Therefore
\[
 \Phi_q^i(w)
 \leq f(\delta)-\delta f'(\delta)
 \leq f(0)=\Phi_p^i(z),
\]
as required.
\end{proof}

\subsection{The largest-root bound and conclusion}

\begin{lemma}[Largest root of the mixed polynomial]
\label{lem:mss-root-bound}
If $A_i\succeq0$, $\sum_iA_i=I_d$, and
$\operatorname{tr}A_i\leq\epsilon$ for every $i$, where $\epsilon>0$,
then the largest root of $\mu[A_1,\ldots,A_m]$ is at most
$(1+\sqrt\epsilon)^2$.
\end{lemma}

\begin{proof}
Let
\[
 p_0(y)=\det\left(\sum_{i=1}^m y_iA_i\right),
 \qquad
 p_j=(1-\partial_{y_j})p_{j-1}\quad(1\leq j\leq m).
\]
The first polynomial is real stable by
Lemma~\ref{lem:mss-stability}, since $\sum_iA_i=I_d$.
Moreover,
\begin{equation}
 \mu[A_1,\ldots,A_m](x)=p_m(x,\ldots,x).
 \label{eq:mss-diagonal-identity}
\end{equation}
To check this identity, substitute $y_i=x+z_i$ before applying the
operators.  Then $\sum_i y_iA_i=xI_d+\sum_i z_iA_i$,
$\partial_{y_i}$ acts as $\partial_{z_i}$, and setting every $z_i=0$
is exactly \eqref{eq:mss-mixed-polynomial}.

Set
\[
 t=\epsilon+\sqrt\epsilon,
 \qquad \delta=1+\sqrt\epsilon,
 \qquad \varphi=\frac\epsilon t=1-\frac1\delta.
\]
At $z^{(0)}=t\mathbf1$, the matrix in $p_0$ is $tI_d$.
Adding nonnegative coordinates adds a positive-semidefinite matrix,
so $z^{(0)}\in\operatorname{Ab}_{p_0}$.
Differentiating the determinant gives, for every $i$,
\[
 \Phi_{p_0}^i(z^{(0)})
 =\operatorname{tr}((tI_d)^{-1}A_i)
 \leq\frac\epsilon t=\varphi.
\]
Inductively define
$z^{(j)}=t\mathbf1+\delta(e_1+\cdots+e_j)$.
Suppose $z^{(j-1)}\in\operatorname{Ab}_{p_{j-1}}$ and all its
barriers are at most $\varphi$.
Because $\varphi+1/\delta=1$, Lemma~\ref{lem:mss-barrier-shift}
applies in direction $j$ and yields
$z^{(j)}\in\operatorname{Ab}_{p_j}$, with all its barriers still
at most $\varphi$.  This proves the induction for every $j$.

At the end, $z^{(m)}=(t+\delta)\mathbf1$ is above the roots of
$p_m$.  By \eqref{eq:mss-diagonal-identity}, the mixed polynomial
is positive for every real $x\geq t+\delta$.
It is real-rooted by Lemma~\ref{lem:mss-expectation}; hence its
largest root is at most
\[
 t+\delta=\epsilon+2\sqrt\epsilon+1
          =(1+\sqrt\epsilon)^2.
\]
\end{proof}

\begin{proof}[Proof of Theorem~\ref{thm:mss}]
Put $A_i=\mathbb E[v_iv_i^*]$.  These matrices are Hermitian
positive semidefinite, and \eqref{eq:mss-hypotheses} says that
$\sum_iA_i=I_d$ and
$\operatorname{tr}A_i=\mathbb E\|v_i\|^2\leq\epsilon$.
Lemma~\ref{lem:mss-root-bound} bounds the largest root of their
mixed polynomial by $(1+\sqrt\epsilon)^2$.
Lemma~\ref{lem:mss-selection} then supplies an outcome of positive
probability whose largest eigenvalue obeys the same bound.
The matrix $\sum_i v_iv_i^*$ is positive semidefinite, so this
largest eigenvalue is its operator norm.  This proves
\eqref{eq:mss-conclusion}.
\end{proof}
